% SPDX-License-Identifier: Apache-2.0
\section{The Rule, and What It Leaves}
\label{sec:e-rule}

\begin{rulebox}[The reduction]
Delete every construct Rust already provides. Rename every construct whose
name collides with a Rust concept. What remains is the library.
\end{rulebox}

The renaming clause does one job in practice. A hardware \term{module} and
a Rust \code{mod} are different things with one name, so the hardware one
becomes a \term{unit} and \code{mod} keeps its meaning as a namespace.

\subsection{The vocabulary that remains}

Four of the library's Rust modules hold everything the language still
needs. The library has others, the netlist, foreign modules, formal
checks and address and register maps, but those are what it does with
a design rather than words a design is written in.

\begin{lstlisting}[caption={The vocabulary that survives the reduction. Four Rust modules and two macros, and everything else is a plain struct, trait or function.},label={lst:vocab}]
crate::types::Transaction   // data that moves between units
crate::types::Tag           // a domain and its policy

crate::comp::Unit         // a unit
crate::comp::Bus            // a bundle of channels
crate::comp::Chan           // one channel, its ends Tx and Rx
crate::comp::Signal         // one wire, its ends Out and In
crate::comp::Clock          // a clock domain, as a type

crate::pipeline::*          // operators that may take a cycle. async,
                            // and a pipeline stage once awaited
crate::funcs::*             // operators that cannot. plain fn

interface!                  // a link's members, a struct per role
config!                     // a build, as one item
\end{lstlisting}

Everything else is a plain Rust struct, trait or function.
Table~\ref{tab:reduction} states the mapping.

\begin{table}[H]
\caption{What each construct becomes. Seventeen of twenty are deleted
outright.}
\label{tab:reduction}
\centering
\footnotesize
\begin{tabular}{@{}ll@{}}
\toprule
Was & Is now \\
\midrule
\code{transaction T}   & \code{struct T} plus \code{\#[derive(Transaction)]} \\
\code{module M}        & \code{struct M}, a unit, plus \code{impl Unit} \\
\code{bus B}           & \code{struct B} plus \code{\#[derive(Bus)]} \\
\code{interface I}     & \code{interface!}, its members declared once \\
\code{modport R}       & \code{struct R}, which \code{interface!} writes \\
\code{channel c: T}    & a field of type \code{Chan<T>} \\
\code{port p}          & an end the unit's \code{run} takes \\
\code{flow f}          & \code{fn f} \\
\code{seq s}           & \code{async fn} \\
\code{pipeline stage pipe} & \code{async fn} and \code{.await} \\
\code{tag T}           & \code{struct T} plus \code{impl Tag}, as a parameter \\
\code{config bind}     & a \code{type} alias \\
\code{domain D}        & a type parameter \\
\code{assert assume cover} & Rust's own macros \\
\bottomrule
\end{tabular}
\end{table}

\subsection{Transactions, interfaces and modports}

A transaction is a struct that implements a marker trait. Nothing else
about it is special.

\begin{lstlisting}[caption={A transaction is a struct and a derive. Nothing else about it is special.},label={lst:transaction}]
#[derive(Transaction)]
pub struct MacRequest {
    pub addr: u32,
    pub count: u8,
}
\end{lstlisting}

An \code{impl} may legally precede the \code{struct} it is for, because
items in a Rust module are not order dependent, and writing
\code{impl Transaction for MacRequest} above the struct would put the kind
before the data the way \code{transaction MacRequest \{ .. \}} used to.
The derive reads the same way and stays idiomatic. It also scales, because
a struct collects several markers before long, and it has somewhere to
grow: the layout type that Section~\ref{sec:e-costs} needs for structural
compatibility is work for this derive to do later.

An interface splits into two kinds of Rust item, and that split fixes a
defect the earlier analysis reported~\cite{unification}. LHDL's
\code{interface} did two jobs at once, a bundle of wires and a socket an
implementation fills, which is why its grammar and its examples never
agreed. Here \code{interface!} declares the members once, and writes a
struct per role whose fields are the ends that role holds: a driving end
for a member it sends on, a reading end for one it takes. The macro
refuses a role that leaves a member out, and two roles that both drive
one member.

\lstinputlisting[caption={An interface declared once, with a struct per role: \code{Initiator}, \code{Target} and \code{Monitor} each hold the ends of the members they use. The probe that answered the question, \code{experiments/rust\_embedding/probe\_interface\_proc.rs}.},label={lst:interface}]{experiments/rust_embedding/probe_interface_proc.rs}

A unit takes a role's struct as a side of its \code{run}, and the
netlist flattens the struct into its members, under their names.
